% GENERATED FILE. Edit canonical lessons, then run npm run generate:books.
% canonical-source-hash: fnv1a64:2df21f8454297211
% canonical-lessons: GU-C287-bhal, GU-C287-paglu, GU-C287-khakhdavvu, GU-C287-santavu, GU-C287-todvu

\chapter{Clues at the Door}
\label{ch:gu-clues-at-the-door}
\hlchaptermodality{287}

\begin{chapteropening}
\textbf{By the end of this chapter:} \emph{I can say what the thief left behind and how he got in.}

Complete the last lesson of chapter 287: I can say what the thief left behind and how he got in.
\end{chapteropening}

\section[\texorpdfstring{\gu{ભાળ} (bhāḷ)}{ભાળ (bhāḷ)}]{\gu{ભાળ} (bhāḷ) — a trace; whereabouts}

\label{lesson:GU-C287-bhal}



\begin{quote}
Before the new one: say the Gujarati for very, then the Gujarati for to hide (something).
\end{quote}

\subsection*{You'll want to know: \gu{ભાળ}}
\textbf{\gu{ભાળ}} — \emph{bhāḷ} — \textquotedblleft{}a trace; whereabouts\textquotedblright{}.

\begin{quote}
\gu{પોલીસ} \gu{ચોરની} \gu{ભાળ} \gu{કાઢતી} \gu{હતી}.
\end{quote}

\begin{quote}
\emph{polīs cornī bhāḷ kāḍhtī hatī.}
\end{quote}

\begin{quote}
The police were trying to trace the thief.
\end{quote}

A feminine noun.

\begin{cousinweb}
\textbf{\gu{ભાળ}} (\emph{bhāḷ}) comes from \emph{*bhāla}, \textquotedblleft{}observation\textquotedblright{}, a form CDIAL reconstructs, and CDIAL lists Gujarati \emph{bhāḷ}, \textquotedblleft{}search, investigation\textquotedblright{} (CDIAL 9473). It is the trail that leads to someone: \textbf{\gu{ભાળ} \gu{કાઢવી}}, to trace a person.
\end{cousinweb}

\begin{grammarlens}[title={What you've built}]
One more word for telling what happened.
\end{grammarlens}

\subsection*{Your turn}
\begin{itemize}
\raggedright
  \item \emph{Say it:} \emph{bhāḷ}
  \item \emph{Say it:} \emph{bhāḷ}, once more
  \item \emph{Say it:} \emph{polīs cornī bhāḷ kāḍhtī hatī.} — The police were trying to trace the thief.
  \item \emph{Recall:} say the Gujarati for very, then the Gujarati for to hide (something), then say \emph{bhāḷ} again
\end{itemize}

\begin{tcolorbox}[breakable,colback=teal!4,colframe=teal!35!black,title={Before you move on}]
What does \gu{ભાળ} mean? (\textquotedblleft{}A trace; whereabouts\textquotedblright{}.) Say it once more.
\end{tcolorbox}
\section[\texorpdfstring{\gu{પગલું} (paglũ)}{પગલું (paglũ)}]{\gu{પગલું} (paglũ) — a footstep}

\label{lesson:GU-C287-paglu}



\begin{quote}
Before the new one: say the Gujarati for to hide (something), then the Gujarati for a trace.
\end{quote}

\subsection*{You'll want to know: \gu{પગલું}}
\textbf{\gu{પગલું}} — \emph{paglũ} — \textquotedblleft{}a footstep\textquotedblright{}.

\begin{quote}
\gu{હવે} \gu{બધું} \gu{આ} \gu{પગલાં} \gu{પર} \gu{આધાર} \gu{રાખે} \gu{છે}.
\end{quote}

\begin{quote}
\emph{have badhũ ā paglā̃ par ādhār rākhe che.}
\end{quote}

\begin{quote}
Now everything rests on these footprints.
\end{quote}

A neuter noun.

\begin{cousinweb}
\textbf{\gu{પગલું}} (\emph{paglũ}) is built on \textbf{\gu{પગ}}, a foot, which CDIAL sets under Sanskrit \emph{padga-} (CDIAL 7766), with the ending \emph{-lũ}. It is a step, and the sound or the print a step leaves: \textbf{\gu{કોઈનાં} \gu{પગલાં}}, someone's footsteps.
\end{cousinweb}

\begin{grammarlens}[title={What you've built}]
One more word for telling what happened.
\end{grammarlens}

\subsection*{Your turn}
\begin{itemize}
\raggedright
  \item \emph{Say it:} \emph{paglũ}
  \item \emph{Say it:} \emph{paglũ}, once more
  \item \emph{Say it:} \emph{have badhũ ā paglā̃ par ādhār rākhe che.} — Now everything rests on these footprints.
  \item \emph{Recall:} say the Gujarati for to hide (something), then the Gujarati for a trace, then say \emph{paglũ} again
\end{itemize}

\begin{tcolorbox}[breakable,colback=teal!4,colframe=teal!35!black,title={Before you move on}]
What does \gu{પગલું} mean? (\textquotedblleft{}A footstep\textquotedblright{}.) Say it once more.
\end{tcolorbox}
\section[\texorpdfstring{\gu{ખખડાવવું} (khakhḍāvvũ)}{ખખડાવવું (khakhḍāvvũ)}]{\gu{ખખડાવવું} (khakhḍāvvũ) — to knock (at a door)}

\label{lesson:GU-C287-khakhdavvu}



\begin{quote}
Before the new one: say the Gujarati for a trace, then the Gujarati for a footstep.
\end{quote}

\subsection*{You'll want to know: \gu{ખખડાવવું}}
\textbf{\gu{ખખડાવવું}} — \emph{khakhḍāvvũ} — \textquotedblleft{}to knock (at a door)\textquotedblright{}.

\begin{quote}
\gu{અડધી} \gu{રાત્રે} \gu{કોઈએ} \gu{બારણું} \gu{ખખડાવ્યું}.
\end{quote}

\begin{quote}
\emph{aḍdhī rātre koīe bārṇũ khakhḍāvyũ.}
\end{quote}

\begin{quote}
In the middle of the night someone knocked at the door.
\end{quote}

Past, with \gu{તેણે} (he or she …): \gu{તેણે} \gu{બારણું} \gu{ખખડાવ્યું}, \emph{teṇe bārṇũ khakhḍāvyũ}.

\begin{cousinweb}
\textbf{\gu{ખખડાવવું}} (\emph{khakhḍāvvũ}) is the causative of \emph{khakhaḍvũ}, to rattle, an imitative word: CDIAL sets Gujarati \emph{khaṛkhaṛvũ}, \textquotedblleft{}to rattle\textquotedblright{}, under Sanskrit \emph{khaṭakhaṭāyatē}, \textquotedblleft{}crackles\textquotedblright{} (CDIAL 3771), and this verb probably belongs with it. To knock is to make the door rattle.
\end{cousinweb}

\begin{grammarlens}[title={What you've built}]
One more word for telling what happened.
\end{grammarlens}

\subsection*{Your turn}
\begin{itemize}
\raggedright
  \item \emph{Say it:} \emph{khakhḍāvvũ}
  \item \emph{Say it:} \emph{khakhḍāvvũ}, once more
  \item \emph{Say it:} \emph{aḍdhī rātre koīe bārṇũ khakhḍāvyũ.} — In the middle of the night someone knocked at the door.
  \item \emph{Recall:} say the Gujarati for a trace, then the Gujarati for a footstep, then say \emph{khakhḍāvvũ} again
\end{itemize}

\begin{tcolorbox}[breakable,colback=teal!4,colframe=teal!35!black,title={Before you move on}]
What does \gu{ખખડાવવું} mean? (\textquotedblleft{}To knock (at a door)\textquotedblright{}.) Say it once more.
\end{tcolorbox}
\section[\texorpdfstring{\gu{સંતાવું} (santāvũ)}{સંતાવું (santāvũ)}]{\gu{સંતાવું} (santāvũ) — to hide oneself}

\label{lesson:GU-C287-santavu}



\begin{quote}
Before the new one: say the Gujarati for a footstep, then the Gujarati for to knock.
\end{quote}

\subsection*{You'll want to know: \gu{સંતાવું}}
\textbf{\gu{સંતાવું}} — \emph{santāvũ} — \textquotedblleft{}to hide oneself\textquotedblright{}.

\begin{quote}
\gu{ચોર} \gu{દાદર} \gu{નીચે} \gu{સંતાઈ} \gu{ગયો}.
\end{quote}

\begin{quote}
\emph{cor dādar nīce santāī gayo.}
\end{quote}

\begin{quote}
The thief hid under the stairs.
\end{quote}

Past (he …): \gu{તે} \gu{ચોર} \gu{સંતાઈ} \gu{ગયો}, \emph{te cor santāī gayo}.

\begin{cousinweb}
goes back to Sanskrit \emph{sáṁtāyatē}, \textquotedblleft{}is stretched out, is covered\textquotedblright{}, and CDIAL lists the Gujarati verb, \textquotedblleft{}to conceal oneself\textquotedblright{} (CDIAL 12887): to be covered over, out of sight. A person hides himself with \textbf{\gu{સંતાવું}}: \textbf{\gu{સંતાઈ} \gu{ગયો}}; a thing is hidden with \textbf{\gu{છુપાવવું}}.
\end{cousinweb}

\begin{grammarlens}[title={What you've built}]
One more word for telling what happened.
\end{grammarlens}

\subsection*{Your turn}
\begin{itemize}
\raggedright
  \item \emph{Say it:} \emph{santāvũ}
  \item \emph{Say it:} \emph{santāvũ}, once more
  \item \emph{Say it:} \emph{cor dādar nīce santāī gayo.} — The thief hid under the stairs.
  \item \emph{Recall:} say the Gujarati for a footstep, then the Gujarati for to knock, then say \emph{santāvũ} again
\end{itemize}

\begin{tcolorbox}[breakable,colback=teal!4,colframe=teal!35!black,title={Before you move on}]
What does \gu{સંતાવું} mean? (\textquotedblleft{}To hide oneself\textquotedblright{}.) Say it once more.
\end{tcolorbox}
\section[\texorpdfstring{\gu{તોડવું} (toḍvũ)}{તોડવું (toḍvũ)}]{\gu{તોડવું} (toḍvũ) — to break (something)}

\label{lesson:GU-C287-todvu}



\begin{quote}
Before the new one: say the Gujarati for to knock, then the Gujarati for to hide oneself.
\end{quote}

\subsection*{You'll want to know: \gu{તોડવું}}
\textbf{\gu{તોડવું}} — \emph{toḍvũ} — \textquotedblleft{}to break (something)\textquotedblright{}.

\begin{quote}
\gu{ચોરે} \gu{તાળું} \gu{તોડ્યું}.
\end{quote}

\begin{quote}
\emph{core tāḷũ toḍyũ.}
\end{quote}

\begin{quote}
The thief broke the lock.
\end{quote}

Past, with \gu{તેણે} (he or she …): \gu{તેણે} \gu{તાળું} \gu{તોડ્યું}, \emph{teṇe tāḷũ toḍyũ}.

\begin{cousinweb}
\textbf{\gu{તોડવું}} (\emph{toḍvũ}) goes back to Sanskrit \emph{trōṭáyati}, \textquotedblleft{}breaks, tears\textquotedblright{}, and CDIAL lists the Gujarati verb (CDIAL 6079). It is the doing partner of \textbf{\gu{તૂટવું}}, to get broken: \textbf{\gu{તાળું} \gu{તૂટ્યું}} says that it broke, \textbf{\gu{ચોરે} \gu{તાળું} \gu{તોડ્યું}} says who broke it.
\end{cousinweb}

\begin{grammarlens}[title={What you've built}]
One more word for telling what happened.
\end{grammarlens}

\subsection*{Your turn}
\begin{itemize}
\raggedright
  \item \emph{Say it:} \emph{toḍvũ}
  \item \emph{Say it:} \emph{toḍvũ}, once more
  \item \emph{Say it:} \emph{core tāḷũ toḍyũ.} — The thief broke the lock.
  \item \emph{Recall:} say the Gujarati for to knock, then the Gujarati for to hide oneself, then say \emph{toḍvũ} again
\end{itemize}

\begin{tcolorbox}[breakable,colback=teal!4,colframe=teal!35!black,title={Before you move on}]
What does \gu{તોડવું} mean? (\textquotedblleft{}To break (something)\textquotedblright{}.) Say it once more.
\end{tcolorbox}
